% !TeX root = ../main.tex
\section{数学における一般性と帰納法、演繹法}

これらは、数学を学習していく上で元来、非常に重要な概念である。数学を学習する意味や意義に通ずる根幹の部分である。

まずは、それらの言葉の意味について確認しておこう。

\begin{definitionbox}{\textbf{一般性}}
    数学における一般性とは、\textbf{限られた具体例にとらわれず、より広い対象や状況に共通する本質的な性質}を指す。
\end{definitionbox}

\begin{definitionbox}{\textbf{帰納法}}
    \textbf{複数の具体的な事例や事実を集めて共通点を見つけ出し、そこから一般的な法則や結論を導き出す論理的推論の方法}を帰納法という。
\end{definitionbox}

\begin{definitionbox}{\textbf{演繹法}}
    \textbf{一般的な法則や結論から具体的な事例や事実を導き出す論理的推論の方法}を演繹法という。
\end{definitionbox}

次に、それらの使い分けについて確認しておこう。


\newpage
